% original_pdf: source/普通高考/2014/2014重庆文.pdf
% page: 1; pdfimages -list found no embedded image objects (the histogram is vector line art)
% grid evidence: tmp/redraw_2014_cq/page300_grid100.png and tmp/redraw_2014_cq/q17_orig_chart.png
% elements: frequency-density y-axis with arrow, score x-axis with arrow and break, five outlined bins,
%            dashed guides at 2a, 3a, 6a, and 7a, ticks 0, 50, 60, 70, 80, 90, 100
% relations: bins [50,60), [60,70), [70,80), [80,90), [90,100) have heights
%            2a, 3a, 7a, 6a, 2a respectively; 20 students and width 10 imply
%            10*(2a+3a+7a+6a+2a)=20, hence a=1/10
% solid segments: bar outlines, x/y axes, tick marks, and the axis-break zigzag
% dashed segments: horizontal guides from the y-axis to the bars at 2a, 3a, 6a, and 7a
% labels: frequency/组距 above-left of y-axis, 2a/3a/6a/7a at guide levels,
%         score ticks below x-axis, 成绩（分） to the right of the final tick
\pgfmathsetmacro{\aheight}{0.1}
\pgfmathsetmacro{\hTwo}{2*\aheight}
\pgfmathsetmacro{\hThree}{3*\aheight}
\pgfmathsetmacro{\hSix}{6*\aheight}
\pgfmathsetmacro{\hSeven}{7*\aheight}
\begin{tikzpicture}[baseline]
\begin{axis}[
    width=8.2cm,
    height=5.4cm,
    axis lines=none,
    clip=false,
    xmin=-12,
    xmax=126,
    ymin=-0.36,
    ymax=0.82,
    ticks=none,
]
    % Dashed guide levels are the frequency-density labels shown in the source.
    \draw[dashed,line width=0.55pt] (axis cs:0,\hTwo) -- (axis cs:100,\hTwo);
    \draw[dashed,line width=0.55pt] (axis cs:0,\hThree) -- (axis cs:70,\hThree);
    \draw[dashed,line width=0.55pt] (axis cs:0,\hSix) -- (axis cs:90,\hSix);
    \draw[dashed,line width=0.55pt] (axis cs:0,\hSeven) -- (axis cs:80,\hSeven);

    \addplot[
        ybar interval,
        draw=black,
        fill=white,
        line width=0.7pt,
    ] coordinates {
        (50,\hTwo)
        (60,\hThree)
        (70,\hSeven)
        (80,\hSix)
        (90,\hTwo)
        (100,0)
    };

    \draw[->,line width=0.8pt] (axis cs:0,0) -- (axis cs:118,0);
    \draw[->,line width=0.8pt] (axis cs:0,0) -- (axis cs:0,0.76);

    % The source omits the interval between 0 and 50, marked by a small zigzag.
    \draw[line width=0.75pt]
      (axis cs:9,0) -- (axis cs:11,0.045) -- (axis cs:13,-0.045) --
      (axis cs:15,0.045) -- (axis cs:17,-0.045) -- (axis cs:19,0);

    \node[anchor=south west,inner sep=0pt,font=\small,align=center]
      at (axis cs:-2,0.70) {\(\dfrac{\text{频率}}{\text{组距}}\)};
    \node[anchor=west,inner sep=0pt,font=\small] at (axis cs:103,0.12) {成绩（分）};

    \draw[line width=0.45pt] (axis cs:0,0) -- (axis cs:0,-0.065);
    \node[anchor=north,inner sep=0pt,font=\small] at (axis cs:0,-0.08) {0};
    \draw[line width=0.45pt] (axis cs:50,0) -- (axis cs:50,-0.065);
    \node[anchor=north,inner sep=0pt,font=\small] at (axis cs:50,-0.08) {50};
    \draw[line width=0.45pt] (axis cs:60,0) -- (axis cs:60,-0.065);
    \node[anchor=north,inner sep=0pt,font=\small] at (axis cs:60,-0.08) {60};
    \draw[line width=0.45pt] (axis cs:70,0) -- (axis cs:70,-0.065);
    \node[anchor=north,inner sep=0pt,font=\small] at (axis cs:70,-0.08) {70};
    \draw[line width=0.45pt] (axis cs:80,0) -- (axis cs:80,-0.065);
    \node[anchor=north,inner sep=0pt,font=\small] at (axis cs:80,-0.08) {80};
    \draw[line width=0.45pt] (axis cs:90,0) -- (axis cs:90,-0.065);
    \node[anchor=north,inner sep=0pt,font=\small] at (axis cs:90,-0.08) {90};
    \draw[line width=0.45pt] (axis cs:100,0) -- (axis cs:100,-0.065);
    \node[anchor=north,inner sep=0pt,font=\small] at (axis cs:100,-0.08) {100};

    \draw[line width=0.45pt] (axis cs:0,\hTwo) -- (axis cs:-2,\hTwo);
    \node[anchor=east,inner sep=0pt,font=\small] at (axis cs:-3,\hTwo) {$2a$};
    \draw[line width=0.45pt] (axis cs:0,\hThree) -- (axis cs:-2,\hThree);
    \node[anchor=east,inner sep=0pt,font=\small] at (axis cs:-3,\hThree) {$3a$};
    \draw[line width=0.45pt] (axis cs:0,\hSix) -- (axis cs:-2,\hSix);
    \node[anchor=east,inner sep=0pt,font=\small] at (axis cs:-3,\hSix) {$6a$};
    \draw[line width=0.45pt] (axis cs:0,\hSeven) -- (axis cs:-2,\hSeven);
    \node[anchor=east,inner sep=0pt,font=\small] at (axis cs:-3,\hSeven) {$7a$};
\end{axis}
\end{tikzpicture}
